\Fidel{TODO
\begin{itemize}
    % \item Justify why we're calculating the covariance or correlation functions.
    \item More detail on non-diagonalisable transfer operators. I think for the sake of writing what we're comfortable with, we can leave it as "assume diagonality", citing the fact that generic choices of process tensors (random SE unitaries) are diagonalisable with probability 1.
    \item Include a generalisation of the transfer operator to higher dimensional SPPs. For example, for an SPP PEPS, the transfer operator itself has an MPO representation. Solving the eigenspectrum then becomes a boundary MPS problem.
    \item Conditional mutual information based definition of approximate Markov order. Leaning towards not including this in the end.
    \item HMM representation and mapping.
\end{itemize}
}

\section{Structure of temporal correlations in 1D SPPs}\label{sec:correlation-structure}
What we'd like to study are correlations between Pauli errors in space and time...

\subsection{Correlation functions}\label{sec:correlation-functions}
We begin by defining correlation functions for SPPs to quantify the statistical dependencies between Pauli operators across space and time. Since SPPs directly encode the joint probability distribution, these correlations can be straightforwardly evaluated via tensor contractions, which later motivates transfer operator formalism.

Let \(f : \mathbb{P}^{(n)} \to \mathbb{R}\) be a scalar function on the set of \(n\)-qubit Pauli operators \(\mathbb{P}^{(n)}\); for example, \(f_X(x) = 1\) if \(x=X\) and 0 otherwise. These functions play a role analogous to observables in quantum theory. Given a random Pauli operator \(x_t \in \mathbb{P}^{(n)}\) at time \(t\), the quantity \(f(x_t)\) defines a real-valued random variable. For two such variables \(f(x_t)\) and \(g(x_{t+\tau})\) evaluated at times \(t\) and \(t+\tau\), the two-point correlation function (covariance) is defined as 
\begin{equation}
    \begin{split}
        C_{f,g}(\tau) 
        &= \mathbb{E}[(f(x_t) - \mathbb{E}[f(x_t)])(g(x_{t+\tau}) - \mathbb{E}[g(x_{t+\tau})])] \\
        &= \mathbb{E}[f(x_t) g(x_{t+\tau})] - \mathbb{E}[f(x_t)] \mathbb{E}[g(x_{t+\tau})],
    \end{split}
\end{equation}
where the expectation \(\mathbb{E}[\cdot]\) is taken with respect to a full joint probability distribution \(\Pr{(x_0,\dots,x_k)}\). At this stage, no assumptions are made regarding properties of the process such as time-homogeneity or stationarity.

In tensor network notation, each function \(f(x_t)\) and \(g(x_{t+\tau})\) corresponds to a real vector \(\mathbf{f}^{x_t}\) and \(\mathbf{g}^{x_{t+\tau}}\) indexed by Pauli labels. The expectation values can then be computed via the full tensor contraction of these vectors with the joint probability tensor \(\mathbf{p}_{x_0\dots x_k}\) (or its MPS representation) and summing over all remaining open indices. Concretely, for single and two-point expectations, we have
\begin{equation}
    \begin{gathered}
        \mathbb{E}[f(x_t)]
        = \sum_{x_0,\dots,x_k} \mathbf{p}_{x_0\dots x_k} 
        \mathbf{f}^{x_t}, \\
        \mathbb{E}[f(x_t) g(x_{t+\tau})] 
        = \sum_{x_0,\dots,x_k} \mathbf{p}_{x_0\dots x_k} 
        \mathbf{f}^{x_t} \mathbf{g}^{x_{t+\tau}},
    \end{gathered}
\end{equation}
respectively.

This construction extends directly to multi-point expectations by inserting the appropriate function vectors. To analyse these correlations more systematically, we now introduce a transfer operator-based framework.

\begin{definition}[Transfer and emission operators]
Consider the MPS representation of an SPP with local tensors \(\mathbf{A}^{\mu_t}{}_{x_t\mu_{t+1}}\) where \(x_t \in \mathbb{P}^{(n)}\). In standard notation, this is equivalently the set of \(\chi \times \chi\) matrices \(A_{x_t}\) (where \(\chi = \dim{\mu}\) is the bond dimension) indexed by the \(x_t\), with elements given by
\begin{equation*}
    (A_{x_t})_{\mu_t \mu_{t+1}} = \mathbf{A}^{\mu_t}{}_{x_t \mu_{t+1}}.
\end{equation*}
Given a scalar function \(f : \mathbb{P}^{(n)} \to \mathbb{R}\), the transfer operator \(T\) and the emission operator \(E_f\) are defined as 
\begin{equation}
    T = \sum_{x_t\in\mathbb{P}^{(n)}}A_{x_t},
    \qquad
    E_f = \sum_{x_t\in\mathbb{P}^{(n)}} f(x_t) A_{x_t}.
\end{equation}
\end{definition}
The transfer operator \(T\) governs how correlations propagate along the virtual (environment) bonds. The observable-inserted transfer operator \(E_f\), which we refer to as the emission operator, encodes how an observable \(f\) couples to these bonds. A tensor graphical representation of these operators is shown in \cref{fig:transfer-emission-covariance}.

We note that the definition of the transfer operator here differs from the conventional quantum MPS transfer operator, which involves both the tensor and its complex conjugate~\cite{schollwockDensitymatrixRenormalizationGroup2005}. The distinction arises because the SPP MPS directly encodes the joint probability distribution over Pauli labels, rather than quantum amplitudes.

We now assume the SPP is \emph{time-homogeneous} (all tensors \(A_{x_t}\) are identical for all \(t\)) and \emph{ergodic}, such that the process reaches a unique stationary distribution on the environment. As a result, the transfer operator \(T\) has a unique leading eigenvalue \(\lambda_1 = 1\) (after normalisation), with corresponding left and right eigenvectors \(\bra{l}\) and \(\ket{r}\) satisfying
\begin{equation}
    \bra{l} T = \bra{l},
    \qquad
    T \ket{r} = \ket{r},
\end{equation}
with normalisation \(\braket{l|r} = 1\). For such stationary processes, the expectation values can be written compactly as
\begin{equation}
    \begin{gathered}
        \mathbb{E}[f(x_t)] = \braket{l|E_f|r} \eqcolon \braket{f}, \\
        \mathbb{E}[f(x_t)g(x_{t+\tau})] = \braket{l|E_f T^{\tau - 1} E_g|r} \eqcolon \braket{fg}.
    \end{gathered}
\end{equation}
To recast the covariance in terms of operators, we introduce the centred emission operator.
\begin{definition}[Centered emission operator]
    The centred emission operator \(\widetilde{E}_f\) corresponding to the function \(f\) is defined as
    \begin{equation}
        \widetilde{E}_f = E_f - \braket{f}\!T,
    \end{equation}
    satisfying \(\braket{l|\widetilde{E}_f|r} = 0\).
\end{definition}
This operator represents the deviation of the observable \(f\) around its stationary mean. Using this, the two-point covariance takes a canonical form,
% \begin{equation}
%     \begin{split}
%         C_{f,g}(\tau) 
%         &= \braket{fg} - \braket{f}\!\!\braket{g}\\
%         &= \braket{fg} - \braket{f}\!\!\braket{g} - \braket{f}\!\!\braket{g} + \braket{f}\!\!\braket{g} \\
%         &= \braket{l|E_f T^{\tau - 1} E_g|r} - \braket{f}\!\!\braket{l|T^\tau E_g|r} - \braket{l|E_f T^\tau|r}\!\!\braket{g} + \braket{f}\!\!\braket{g}\!\!\braket{l|T^{\tau+1}|r}\\
%         &= \braket{l|(E_f - \braket{f}\!T)T^{\tau-1}(E_g - \braket{g}\!T)|r} \\
%         &= \braket{l|\widetilde{E}_f T^{\tau - 1} \widetilde{E}_g|r}.
%     \end{split}
% \end{equation}
\begin{equation}\label{eq:covariance_operator}
    C_{f,g}(\tau) = \braket{l|\widetilde{E}_f T^{\tau - 1} \widetilde{E}_g|r}.
\end{equation}
The covariance is thus described entirely in terms of the transfer and emission operators. It follows that the decay of temporal correlations is governed by the exponent of the transfer operator \(T\), and hence by its eigenvalue spectrum. Finally, this formulation extends naturally to multi-point correlation functions. For an ordered time sequence \(t_1 < t_2 < \dots < t_m\), the \(m\)-point correlator is
\begin{equation}
    C_{f_1,\dots,f_m} = \braket{l|\widetilde{E}_{f_1} T^{\Delta t_1} \widetilde{E}_{f_2} T^{\Delta t_2} \dots T^{\Delta t_{m-1}}\widetilde{E}_{f_m}|r},
\end{equation}
where \(\Delta t_j = t_{j+1} - t_j - 1\). We now turn to the spectral properties of the transfer operator to characterise the decay of these correlations over time.

\begin{figure}[t]
    \centering
    \begin{subfigure}{0.33\textwidth}\captionsetup{skip=-1mm}
        \centering
        \[
        \valignboxnopad{\begin{tikzpicture}
	\begin{pgfonlayer}{nodelayer}
		\node [style=smallboxyellow] (0) at (0, 0) {$T$};
		\node [style=none] (1) at (0.75, 0) {};
		\node [style=none] (2) at (-0.75, 0) {};
	\end{pgfonlayer}
	\begin{pgfonlayer}{edgelayer}
		\draw [style=fusedwire] (0) to (1.center);
		\draw [style=fusedwire] (2.center) to (0);
	\end{pgfonlayer}
\end{tikzpicture}
} 
        \!=\;
        \sum_x \valignboxnopad{\begin{tikzpicture}
	\begin{pgfonlayer}{nodelayer}
		\node [style=smallcircleorange] (0) at (0, 0) {\(\mathbf{A}\)};
		\node [style=none] (1) at (0.75, 0) {};
		\node [style=none] (2) at (-0.75, 0) {};
		\node [style=none] (3) at (0, -0.675) {};
		\node [style=none] (4) at (0, -0.8) {\({}_x\)};
	\end{pgfonlayer}
	\begin{pgfonlayer}{edgelayer}
		\draw [style=fusedwire] (0) to (1.center);
		\draw [style=fusedwire] (2.center) to (0);
		\draw (3.center) to (0);
	\end{pgfonlayer}
\end{tikzpicture}
}
        \]
        \caption{}
    \end{subfigure}
    \begin{subfigure}{0.33\textwidth}\captionsetup{skip=-1mm}
        \centering
        \[
        \valignboxnopad{\begin{tikzpicture}
	\begin{pgfonlayer}{nodelayer}
		\node [style=smallboxblue] (0) at (0, 0) {\(E_f\)};
		\node [style=none] (1) at (0.75, 0) {};
		\node [style=none] (2) at (-0.75, 0) {};
	\end{pgfonlayer}
	\begin{pgfonlayer}{edgelayer}
		\draw [style=fusedwire] (1.center) to (0);
		\draw [style=fusedwire] (2.center) to (0);
	\end{pgfonlayer}
\end{tikzpicture}
} 
        \!=
        \valignboxnopad{\begin{tikzpicture}
	\begin{pgfonlayer}{nodelayer}
		\node [style=smallcircleorange] (0) at (0, 0) {\(\mathbf{A}\)};
		\node [style=none] (1) at (0.75, 0) {};
		\node [style=none] (2) at (-0.75, 0) {};
		\node [style=downtriangleblue] (3) at (0, -0.825) {\(\mathbf{f}\)};
	\end{pgfonlayer}
	\begin{pgfonlayer}{edgelayer}
		\draw [style=fusedwire] (0) to (1.center);
		\draw [style=fusedwire] (2.center) to (0);
		\draw (0) to (3);
	\end{pgfonlayer}
\end{tikzpicture}
}
        \]
        \caption{}
    \end{subfigure}
    \begin{subfigure}{0.7\textwidth}\captionsetup{skip=-2mm}
        \centering
        \[
        C_{f,g}(\tau) \;=\;
        \raisebox{-1mm}{\valignboxnopad{\begin{tikzpicture}
	\begin{pgfonlayer}{nodelayer}
		\node [style=smallboxblue] (0) at (0.75, 0) {\(\widetilde{E}_f\)};
		\node [style=smallboxyellow] (1) at (1.75, 0) {\(T\)};
		\node [style=smallboxyellow] (2) at (3.25, 0) {\(T\)};
		\node [style=lefttrianglesmall] (3) at (0, 0) {\(l\)};
		\node [style=smallboxblue] (4) at (4.25, 0) {\(\widetilde{E}_g\)};
		\node [style=righttrianglesmall] (5) at (5, 0) {\(r\)};
		\node [style=none] (6) at (2.75, 0) {};
		\node [style=none] (7) at (2.25, 0) {};
		\node [style=none] (8) at (1.375, -0.475) {};
		\node [style=none] (9) at (3.625, -0.475) {};
		\node [style=none] (10) at (2.5, -0.875) {\({}_{\tau - 1}\)};
	\end{pgfonlayer}
	\begin{pgfonlayer}{edgelayer}
		\draw [style=fusedwire] (3) to (0);
		\draw [style=fusedwire] (0) to (1);
		\draw [style=fusedwire] (4) to (2);
		\draw [style=fusedwire] (4) to (5);
		\draw [style=fusedwire] (2) to (6.center);
		\draw [style=fusedwire] (1) to (7.center);
		\draw [style=fuseddotted] (7.center) to (6.center);
		\draw [style=braceunder, in=180, out=0] (8.center) to (9.center);
	\end{pgfonlayer}
\end{tikzpicture}
}}
        \]
        \caption{}
    \end{subfigure}
    \caption{}
    \label{fig:transfer-emission-covariance}
\end{figure}

\subsection{Spectral analysis of SPP transfer operators}\label{sec:spectral-analysis}
For clarity, we begin by assuming that the transfer operator \(T\) is diagonalisable. Its spectral decomposition is given by
\begin{equation}
    T = \sum_{i=1}^{\chi} \lambda_i \ketbra{r_i}{l_i},
\end{equation}
where the eigenvalues are ordered by their magnitudes \(|\lambda_1| \geq |\lambda_2| \geq \dots \geq |\lambda_{\chi}|\), \(\braket{l_i|r_j} = \delta_{ij}\), and normalisation \(\lambda_1 = 1\). Inserting this into the covariance function from \cref{eq:covariance_operator} yields
\begin{equation}
    C_{f,g}(\tau) = \sum_{i=2}^{\chi} \lambda_i^{\tau - 1} \! \braket{l|\widetilde{E}_f|r_i} \! \braket{l_i|\widetilde{E}_g|r}.
\end{equation}
Only the terms \(i \geq 2\) contribute since the leading eigenvectors correspond to the stationary distribution, for which \(\braket{l_1|\widetilde{E}_f|r_1} = \braket{l_1|\widetilde{E}_g|r_1} = 0\). This expression explicitly shows that the decay of correlations is a sum of exponentially decaying modes.

The asymptotic, long-time behaviour is dominated by the subleading eigenvalue with the largest magnitude, which we denote as \(\lambda_* = |\lambda_2|\). From this, we define two key quantities:
\begin{enumerate}
    \item \textbf{The spectral gap} (\(\Delta\)): Defined as 
    \begin{equation}
        \Delta = 1 - \lambda_*,
    \end{equation}
    the spectral gap quantifies the rate of convergence to the stationary distribution. A large gap (\(\lambda_* \ll 1\)) implies correlations decay rapidly, indicating a short memory time. Conversely, a small gap (\(\lambda_* \approx 1\)) signifies long-lived temporal correlations.
    \item \textbf{The correlation length} (\(\xi\)): The characteristic timescale of the decay is given by the correlation length,
    \begin{equation}
        \xi = -\frac{1}{\ln \lambda_*}.
    \end{equation}
    This means the asymptotic decay is governed by this timescale, with \(|C_{f,g}(\tau)| \sim e^{-(\tau - 1)/\xi}\) for sufficiently large \(\tau\).
\end{enumerate}
The eigenspectrum of the transfer operator provides a comprehensive characterisation of how environmental memory is propagated and dissipated. These spectral properties remain a useful descriptor even when relaxing assumptions such as diagonalisability and ergodicity. For instance, complex eigenvalues indicate oscillatory correlation behaviour, while degenerate leading eigenvalues imply non-ergodic dynamics with multiple stationary states.

A more formal bound on the covariance can be established. Let \(S \coloneq \{i \geq 2: \braket{l|\widetilde{E}_f \ketbra{r_i}{l_i} \widetilde{E}_g|r} \neq 0\}\) and \(k \coloneq |S| \leq \chi - 1\). In other words, \(S\) is the set of eigenmodes that contribute non-trivially to the correlation function for the given observables \(f\) and \(g\), and \(k\) is the number of such modes. For a diagonalisable \(T\) working in the canonical gauge \(\braket{l|r} = 1\) and with any submultiplicative norm \(\|\cdot\|\), we obtain the bound,

\begin{equation}
    |C_{f,g}(\tau)| \leq \sum_{i\in S} |\lambda_i|^{\tau - 1} \| \widetilde{E}_f \| \| \widetilde{E}_g \| \leq k\|\widetilde{E}_f \| \| \widetilde{E}_g \| \lambda_*^{\tau - 1},
\end{equation}
If only a few modes matter, \(k\) may be much smaller than \(\chi - 1\). Finally, while \(T=\sum_x A_x\) is generally non-Hermitian (and may be non-normal), generic choices of process tensors---such as those generated by Haar-random system-environment unitaries---are almost surely diagonalisable. Next, we discuss the connection between SPP MPS representations and hidden Markov models.

\subsection{Hidden Markov model representations of SPPs}
An SPP MPS, together with its associated transfer operator, naturally admits a hidden state interpretation. In this view, the virtual bond indices serve as latent variables whose dynamics drive the observed Pauli sequence. We make this connection explicit by linking SPP MPS representations to finite-state hidden Markov models (HMMs). Concretely, under appropriate positivity conditions and normalisation conditions, an SPP MPS of bond dimension \(\chi\) is isomorphic to an \emph{edge-emitting} HMM with \(\chi\) hidden states. This equivalence provides a bridge to the extensive HMM literature, enabling efficient hidden state-based methods for sampling and modelling SPPs---tools we exploit in \cref{sec:temporalstormmodel,sec:qcamodel}.

Let \((S_t)_{t\geq 0}\) be a Markov chain on a finite state space \(\mathcal{S}=\{1,\dots,\chi\}\), and let \((X_t)_{t\geq 0}\) be a sequence of observations from the Pauli alphabet \(\mathbb{P}^{(n)}\). An edge-emitting HMM is fully specified by an initial distribution \(\pi_i\) over the hidden states and a set of transition-emission kernels \((K_x)_{ij}\). The initial distribution satisfies
\begin{equation}
    \pi_i \coloneq \Pr(S_0=i), 
    \quad 
    \pi_i \ge 0,\quad \sum_{i\in\mathcal S}\pi_i = 1.
\end{equation}
The transition-emission kernels \((K_x)_{ij}\) describe the joint probability of transitioning from hidden state \(i\) to \(j \in \mathcal{S}\) while emitting observation \(x \in \mathbb{P}^{(n)}\), defined as
\begin{equation}
    (K_x)_{ij} \coloneq \Pr{(S_{t+1}=j, X_t = x | S_t = i)},
\end{equation}
whose elements are non-negative and appropriately normalised as
\begin{equation}\label{eq:hmm-kernel-conditions}
    (K_x)_{ij} \geq 0 \quad \forall i,j,x,
    \quad
    \sum_{x,j} (K_x)_{ij} = 1 \quad \forall i.
\end{equation}
Summing over emitted symbols yields the hidden state \emph{transition matrix},
% Equivalently, the kernels \((K_x)_{ij}\) can be re-expressed in terms of separate transition and emission probabilities. For example, the transition matrix can be recovered from \((K_x)_{ij}\) as
\begin{equation}
    (T_{\mathrm{HMM}})_{ij} = \sum_{x} (K_x)_{ij},
\end{equation}
which is \emph{row-stochastic} by \cref{eq:hmm-kernel-conditions}, satisfying \(\sum_j (T_{\mathrm{HMM}})_{ij} = 1\) for all \(i\).

The probability of observing a Pauli sequence \(x_{0:k}\) is given by a marginalisation over all possible hidden state sequences,
\begin{equation}
    \Pr(x_{0:k}) = \sum_{s_0,\dots,s_{k+1}} \pi_{s_0} \prod_{t=0}^k (K_{x_t})_{s_t s_{t+1}}.
\end{equation}
This expression is in direct analogy with the SPP MPS probabilities in \cref{eq:spp-mps}. We now propose sufficient conditions under which an SPP MPS is equivalent to such an HMM.
\begin{proposition}[Sufficient conditions for SPP MPS-HMM equivalence]
    Let a time-homogeneous SPP be specified by an MPS \(\{A_x\}_{x\in\mathbb P^{(n)}}\) of bond dimension \(\chi\). If there exists a representation (gauge) in which \(A_x\) satisfies
    \begin{equation}
        (A_x)_{ij}\ge 0 \ \ \forall\, i,j,x, \quad
        \sum_{x\in\mathbb{P}^{(n)}}\sum_{j=1}^{\chi} (A_x)_{ij}=1 \ \ \forall\, i,
    \end{equation}
    then the process admits an edge-emitting Hidden Markov model (HMM) realisation with \(\chi\) hidden states under the identification
    \begin{equation*}
        (K_x)_{ij}\equiv(A_x)_{ij}.
    \end{equation*}
    In particular, the SPP transfer matrix \(T = \sum_x A_x\) coincides with the HMM transition matrix \(T_{\mathrm{HMM}} = \sum_x K_x\).
\end{proposition}
These conditions are sufficient but not necessary; a generic SPP MPS need not admit a nonnegative, row-stochastic representation of the same bond dimension. However, any non-Markovian process can be rendered Markovian by enlarging the (in general, potentially unbounded) state space \cite{milzQuantumStochasticProcesses2021}. Constructing such an embedding optimally is model dependent and generally nontrivial~\cite{vidyasagarCompleteRealizationProblem2011}. Since we are primarily interested in utilising HMMs as practical generative models for correlated noise, we do not pursue minimal HMM realisations here.

A practical advantage of HMM realisations is access to efficient procedures for inference, learning, and sampling. Importantly, HMMs enable simple sequential sampling of Pauli trajectories based on the transition-emission updates, making them well suited to large-scale Monte Carlo studies of QEC under correlated noise. This capability underpins the numerical analyses in subsequent sections.



\subsection{Transfer operators for higher-dimensional SPPs}
Before concluding this section, we briefly indicate how the transfer operator framework extends to higher-dimensional SPP tensor networks. In the one-dimensional case, all spatial structure is compressed into a single index, resulting in a matrix transfer operator. For higher dimensional SPPs, the transfer operator itself may be represented by a tensor network.

Concretely, consider the 2D SPP PEPS illustrated in \cref{fig:spp-peps}, whose local tensors are given by
\begin{equation*}
    \mathbf{a}_{(i, j)}^{\nu^{i-1}_j\mu_j^i}{}_{x_j^i\mu_{j+1}^i\nu_j^i},
\end{equation*}
where \(i\) and \(j\) label spatial and temporal positions, respectively.
We define the \emph{single-site transfer tensor} \(\mathbf{r}_{(i,j)}\) by summing over the local Pauli index \(x_j^i\). That is, for each site \((i,j)\), we obtain
\begin{equation}
    \mathbf{r}_{(i, j)}^{\nu^{i-1}_j\mu_j^i}{}_{\mu_{j+1}^i\nu_j^i} = 
    \sum_{x_j^i \in \mathbb{P}} \mathbf{a}_{(i, j)}^{\nu^{i-1}_j\mu_j^i}{}_{x_j^i\mu_{j+1}^i\nu_j^i}.
\end{equation}
The \emph{row transfer operator} \(\mathbf{R}_{(j)}\) is then the contraction of these tensors along the spatial bonds \(\nu\) across a single row (time slice) \(j\). Writing \(\vec{\mu}_j \coloneq (\mu_j^0, \dots, \mu_j^{n-1})\) for the collection of temporal bonds (and similarly \(\vec{\mu}_{j+1}\)), we have
\begin{equation}
    \mathbf{R}_{(j)}^{\vec{\mu}_j}{}_{\vec{\mu}_{j+1}} = 
    \sum_{\nu^0_j,\dots,\nu^{n-1}_j}
    \mathbf{r}_{(0, j)}^{\mu_j^0}{}_{\mu_{j+1}^0\nu_j^0}
    \mathbf{r}_{(1, j)}^{\nu^0_j\mu_j^1}{}_{\mu_{j+1}^1\nu_j^1}
    \dots
    \mathbf{r}_{(n-1, j)}^{\nu^{n-2}_j\mu_j^{n-1}}{}_{\mu_{j+1}^{n-1}\nu_j^{n-1}},
\end{equation}
which is precisely an MPO along space.

At a high level, the previous analyses of \cref{sec:correlation-functions,sec:spectral-analysis} may be generalised to higher-dimensional SPPs by replacing the matrix transfer operator \(T\) with an appropriate tensor network realisation (such as the MPO \(\mathbf{R}_{(j)}\) for 2D SPP PEPS).